% «Más ejemplos resueltos» · Razonamiento lógico (libros 1 y 2). Salen solos antes de «Así lo preguntan en el examen».
% --- Libro 1 · Cálculo mental ---
\masejemplosde{pu-cm-1}{
\ejx{El saldo de la tarjeta}{Una tarjeta de transporte tiene un saldo de \$3\,200. Se hacen tres viajes de \$2\,950 (la tarjeta permite quedar debiendo) y luego se recargan \$10\,000. ¿Qué saldo queda?}{$3\,200 - 3 \cdot 2\,950 = 3\,200 - 8\,850 = -5\,650$. Con la recarga: $-5\,650 + 10\,000 = 4\,350$.}
\ejx{Una expresión con trampa}{Calcula $\;-4^2 - (-4)^2 + (-1)^{7} \cdot (2 - 5)$.}{$-4^2 = -16$ y $(-4)^2 = 16$; $(-1)^7 = -1$ y $2 - 5 = -3$, cuyo producto es $3$. Resultado: $-16 - 16 + 3 = -29$.}}
\masejemplosde{pu-cm-2}{
\ejx{Lo que falta del camino}{Un ciclista ha recorrido $\tfrac{5}{8}$ de una ruta y le faltan 27 km. ¿Cuánto mide la ruta?}{Le faltan $\tfrac38$, que son 27 km: un octavo son 9 km y la ruta mide $8 \cdot 9 = 72$ km.}
\ejx{¿Cuál es mayor?}{Ordena de menor a mayor $\tfrac{3}{4}$, $\tfrac{7}{10}$ y $\tfrac{11}{15}$.}{Con denominador 60: $\tfrac{45}{60}$, $\tfrac{42}{60}$ y $\tfrac{44}{60}$. Orden: $\tfrac{7}{10} < \tfrac{11}{15} < \tfrac34$.}}
\masejemplosde{pu-cm-3}{
\ejx{La cifra lejana}{¿Qué cifra ocupa el lugar 75 después del punto en $\tfrac{5}{11}$?}{$\tfrac{5}{11} = 0.\overline{45}$: los lugares impares tienen 4 y los pares 5. El 75 es impar: 4.}
\ejx{Tres formatos}{En un grupo, $\tfrac{2}{5}$ prefieren fútbol, el $0.35$ baloncesto y el resto, 20 personas, natación. ¿Cuántos son?}{$\tfrac25 = 40\,\%$ y $0.35 = 35\,\%$: suman $75\,\%$. El $25\,\%$ restante son 20 personas: el grupo tiene 80.}}
\masejemplosde{pu-cm-4}{
\ejx{Misma base}{Simplifica $\dfrac{9^4 \cdot 27}{3^9}$.}{$9^4 = 3^8$ y $27 = 3^3$: $\tfrac{3^{11}}{3^9} = 3^2 = 9$.}
\ejx{Estimar una raíz}{¿Cuál entero está más cerca de $\sqrt{300}$?}{$17^2 = 289$ y $18^2 = 324$. Como 300 está a 11 de 289 y a 24 de 324, el más cercano es 17.}}
% --- Aritmética y promedios ---
\masejemplosde{pu-ar-1}{
\ejx{El ascensor de fracciones}{De un tanque lleno se saca $\tfrac13$; de lo que queda, $\tfrac14$; y de lo que queda, $\tfrac15$. Quedan 120 litros. ¿Cuál es su capacidad?}{Queda $\tfrac23 \cdot \tfrac34 \cdot \tfrac45 = \tfrac{24}{60} = \tfrac25$. Si $\tfrac25$ son 120, un quinto son 60 y el tanque tiene 300 litros.}
\ejx{Tres operaciones al revés}{Un número se triplica, se le restan 9, el resultado se divide entre 4 y da 6. ¿Cuál era?}{Al revés: $6 \cdot 4 = 24$; $24 + 9 = 33$; $33 \div 3 = 11$.}}
\masejemplosde{pu-ar-2}{
\ejx{Sale el mayor}{El promedio de edad de 6 jugadores es 24 años. Si sale el de 34 años, ¿cuál es el promedio de los otros cinco?}{Total: $6 \cdot 24 = 144$. Sin él: $110$. Promedio: $110 \div 5 = 22$.}
\ejx{Dos secciones}{La sección A (20 estudiantes) promedió 3.4 y la B (30 estudiantes) 3.9. ¿Cuál es el promedio de las dos juntas?}{$\tfrac{20 \cdot 3.4 + 30 \cdot 3.9}{50} = \tfrac{68 + 117}{50} = 3.7$.}}
\masejemplosde{pu-ar-3}{
\ejx{Una resta con huecos}{Completa $\;7\,\square\,2 - \square\,6\,5 = 4\,5\,7$.}{Unidades: $12 - 5 = 7$ (se pide prestado). Decenas: $\square - 1 - 6$ termina en 5: $\square = 2$ y se pide prestado ($12 - 1 - 6 = 5$). Centenas: $7 - 1 - \square = 4$: $\square = 2$. Es $722 - 265 = 457$.}
\ejx{La suma de un número y su reverso}{Un número de dos cifras sumado con el que resulta de invertir sus cifras da 99, y la cifra de las decenas es el doble de la de las unidades. ¿Cuál es?}{$11(a + b) = 99$: $a + b = 9$. Con $a = 2b$: $b = 3$, $a = 6$. El número es 63.}}
\masejemplosde{pu-ar-4}{
\ejx{Las llaves}{Un llavero tiene 10 llaves y solo una abre la puerta. Probándolas sin repetir, ¿cuántos intentos se necesitan, como máximo, para saber con certeza cuál es?}{En el peor caso fallan 9 seguidas; la décima ya no hay que probarla, porque es la única que queda. Bastan 9 intentos.}
\ejx{Billetes distintos}{Cuatro amigos reúnen \$50\,000 entre todos, cada uno aporta un número entero de miles y todos aportan cantidades distintas, mínimo \$5\,000. ¿Cuánto puede aportar, como máximo, uno de ellos?}{Los otros tres, lo mínimo y distinto: $5 + 6 + 7 = 18$ mil. El cuarto puede aportar \$32\,000.}}
% --- Operadores ---
\masejemplosde{pu-op-1}{
\ejx{Dos operadores}{Si $a \oplus b = a + 2b$ y $a \otimes b = ab - 1$, calcula $(2 \oplus 3) \otimes 2$.}{$2 \oplus 3 = 2 + 6 = 8$; $8 \otimes 2 = 16 - 1 = 15$.}
\ejx{Un operador con fracción}{Si $m \triangle n = \dfrac{m - n}{m + n}$, calcula $7 \triangle 3$ y $3 \triangle 7$.}{$\tfrac{4}{10} = \tfrac25$ y $\tfrac{-4}{10} = -\tfrac25$: al intercambiar, cambia el signo.}}
\masejemplosde{pu-op-2}{
\ejx{Despejar adentro}{Si $a \star b = a^2 + b$ y $3 \star x = x \star 3$ con $x \ne 3$, halla $x$.}{$9 + x = x^2 + 3$, es decir $x^2 - x - 6 = 0$: $(x - 3)(x + 2) = 0$. Como $x \ne 3$, $x = -2$.}
\ejx{Regla recursiva}{Si $g(1) = 5$ y $g(n + 1) = g(n) + 2n$, calcula $g(5)$.}{$g(2) = 7$, $g(3) = 11$, $g(4) = 17$, $g(5) = 25$.}}
% --- Proporcionalidad ---
\masejemplosde{pu-pr-1}{
\ejx{Receta ampliada}{Una receta para 6 personas lleva 450 g de arroz y 300 g de pollo. ¿Cuánto de cada uno se necesita para 10 personas?}{Por persona: 75 g de arroz y 50 g de pollo. Para 10: 750 g y 500 g.}
\ejx{Tres hermanos}{Tres hermanos se reparten 72 hectáreas en razón $1 : 3 : 5$. ¿Cuántas hectáreas más recibe el mayor que el menor?}{9 partes de 8 hectáreas: reciben 8, 24 y 40. La diferencia es 32.}}
\masejemplosde{pu-pr-2}{
\ejx{El grupo que crece}{Un trabajo que 9 personas hacen en 20 días debe terminarse en 12 días. ¿Cuántas personas más se necesitan?}{Trabajo: $9 \cdot 20 = 180$ persona-días. En 12 días: $180 \div 12 = 15$ personas. Faltan 6.}
\ejx{Compuesta}{Cinco impresoras imprimen 3\,000 afiches en 4 horas. ¿Cuántos afiches imprimen 8 impresoras en 3 horas?}{Cada impresora hace $\tfrac{3\,000}{5 \cdot 4} = 150$ por hora. $8 \cdot 3 \cdot 150 = 3\,600$.}}
\masejemplosde{pu-pr-3}{
\ejx{Ventaja en la carrera}{Un atleta corre a 6 m/s y otro a 4 m/s. Si el lento sale con 100 m de ventaja, ¿en cuánto tiempo lo alcanza el rápido?}{Le descuenta 2 m cada segundo: $100 \div 2 = 50$ s.}
\ejx{Trabajan juntos y uno se va}{Ana digita un informe en 6 horas y Luis en 3. Trabajan juntos 1 hora y Luis se va. ¿Cuánto tarda Ana en terminar sola?}{Juntos hacen $\tfrac16 + \tfrac13 = \tfrac12$ por hora: tras 1 hora falta la mitad. Ana tarda $\tfrac12 \div \tfrac16 = 3$ horas.}}
\masejemplosde{pu-pr-4}{
\ejx{Aprobados}{En un examen aprobaron 36 de 48 estudiantes. ¿Qué porcentaje reprobó?}{Reprobaron 12 de 48: $\tfrac14 = 25\,\%$.}
\ejx{Del porcentaje al total}{El $18\,\%$ de los libros de una biblioteca son de ciencias: 540 libros. ¿Cuántos libros hay en total?}{El $1\,\%$ son $540 \div 18 = 30$ libros; el total, 3\,000.}}
\masejemplosde{pu-pr-5}{
\ejx{Tres años de aumentos}{Un arriendo de \$1\,000\,000 sube $10\,\%$ cada año. ¿Cuánto vale después de 3 años?}{$1\,000\,000 \cdot 1.1^3 = 1\,331\,000$ (no \$1\,300\,000: cada aumento se calcula sobre el valor ya aumentado).}
\ejx{Descuento y recargo}{Un computador de \$2\,000\,000 tiene $15\,\%$ de descuento, y al precio rebajado se le suma el IVA del $19\,\%$. ¿Cuánto se paga?}{$2\,000\,000 \cdot 0.85 = 1\,700\,000$; $1\,700\,000 \cdot 1.19 = 2\,023\,000$.}}
% --- Ecuaciones ---
\masejemplosde{pu-ec-1}{
\ejx{La cuerda cortada}{Una cuerda de 90 cm se corta en tres trozos: el segundo mide el doble del primero y el tercero 10 cm más que el segundo. ¿Cuánto mide cada uno?}{$x + 2x + (2x + 10) = 90$: $5x = 80$, $x = 16$. Miden 16, 32 y 42 cm.}
\ejx{Fracciones en la ecuación}{La tercera parte de un número más su cuarta parte es 21. ¿Cuál es el número?}{$\tfrac{x}{3} + \tfrac{x}{4} = \tfrac{7x}{12} = 21$: $x = 36$.}}
\masejemplosde{pu-ec-2}{
\ejx{Edades con suma}{Hoy las edades de Julia y su abuelo suman 80 años. Hace 10 años, el abuelo tenía 5 veces la edad de Julia. ¿Cuántos años tiene Julia?}{Hace 10 años sumaban 60 y estaban en razón $1 : 5$: Julia tenía 10 y el abuelo 50. Hoy Julia tiene 20.}
\ejx{Billetes}{En una caja hay 40 billetes de \$2\,000 y \$5\,000 que suman \$125\,000. ¿Cuántos hay de \$5\,000?}{Si todos fueran de 2\,000: \$80\,000. Faltan \$45\,000, y cada billete de 5\,000 aporta 3\,000 más: 15 billetes.}}
\masejemplosde{pu-ec-3}{
\ejx{Dos compras}{Tres empanadas y dos jugos cuestan \$13\,000; una empanada y dos jugos, \$7\,000. ¿Cuánto cuesta un jugo?}{Restando: dos empanadas cuestan \$6\,000, cada una \$3\,000. Entonces dos jugos cuestan \$4\,000: un jugo, \$2\,000.}
\ejx{Presupuesto}{Con \$50\,000 se quieren comprar cuadernos de \$4\,500 después de pagar \$8\,000 de transporte. ¿Cuántos cuadernos, como máximo?}{$8\,000 + 4\,500n \le 50\,000$: $n \le 9.3$. Máximo 9 cuadernos.}}
% --- Teoría de números ---
\masejemplosde{pu-nu-1}{
\ejx{La cifra que falta}{¿Qué cifras puede ser $d$ para que $\overline{2d6}$ sea divisible entre 4?}{Las dos últimas, $\overline{d6}$, deben ser múltiplo de 4: 16, 36, 56, 76, 96. $d$ puede ser 1, 3, 5, 7 o 9.}
\ejx{Divisores de un cuadrado}{¿Cuántos divisores tiene 144?}{$144 = 2^4 \cdot 3^2$: $(4 + 1)(2 + 1) = 15$. Es impar, como en todo cuadrado perfecto.}}
\masejemplosde{pu-nu-2}{
\ejx{Cajas iguales}{Se quieren empacar 84 manzanas y 126 naranjas en cajas iguales, sin mezclar frutas y con el mayor número posible de frutas por caja. ¿Cuántas cajas se necesitan?}{MCD$(84, 126) = 42$ frutas por caja: $2 + 3 = 5$ cajas.}
\ejx{Tres semáforos}{Tres semáforos cambian a verde cada 40, 60 y 90 segundos. Si coinciden a las 7:00:00, ¿a qué hora vuelven a coincidir?}{MCM$(40, 60, 90) = 360$ s $= 6$ min: a las 7:06:00.}}
\masejemplosde{pu-nu-3}{
\ejx{Última cifra de un producto}{¿En qué cifra termina $7^{50} \cdot 4^{31}$?}{$7^{50}$: $50 = 4 \cdot 12 + 2$, termina en 9. $4^{31}$: exponente impar, termina en 4. $9 \cdot 4 = 36$: termina en 6.}
\ejx{El cumpleaños}{Si un cumpleaños cae en viernes, ¿en qué día cae el del año siguiente, si no hay 29 de febrero en medio?}{$365 = 52 \cdot 7 + 1$: avanza un día. Cae en sábado.}}
% --- Geometría ---
\masejemplosde{pu-ge-1}{
\ejx{Ángulos en razón}{Los tres ángulos de un triángulo están en razón $2 : 3 : 4$. ¿Cuánto mide el mayor?}{9 partes de $20^\circ$: $40^\circ$, $60^\circ$ y $80^\circ$. El mayor, $80^\circ$.}
\ejx{El reloj}{¿Qué ángulo forman las manecillas a las 4:00? ¿Y a las 4:30?}{A las 4:00, cuatro divisiones de $30^\circ$: $120^\circ$. A las 4:30 el minutero está en $180^\circ$ y el horario en $135^\circ$: $45^\circ$.}}
\masejemplosde{pu-ge-2}{
\ejx{El televisor}{Un televisor tiene una pantalla de 80 cm de ancho y 60 cm de alto. ¿Cuánto mide su diagonal?}{Terna $3$-$4$-$5$ multiplicada por 20: 100 cm.}
\ejx{El poste y el cable}{Un cable de 25 m sostiene un poste desde su punta hasta un punto del suelo a 7 m de la base. ¿Qué altura tiene el poste?}{Terna $7$-$24$-$25$: 24 m.}}
\masejemplosde{pu-ge-3}{
\ejx{El plano del apartamento}{En un plano a escala $1 : 50$, una habitación mide $8 \text{ cm} \times 6 \text{ cm}$. ¿Cuál es su área real?}{Medidas reales: 4 m y 3 m. Área: 12 m$^2$.}
\ejx{Dos triángulos}{Dos triángulos semejantes tienen lados correspondientes de 6 y 9 cm. Si el área del pequeño es 20 cm$^2$, ¿cuál es la del grande?}{$k = \tfrac32$, $k^2 = \tfrac94$: $20 \cdot \tfrac94 = 45$ cm$^2$.}}
\masejemplosde{pu-ge-4}{
\ejx{La ventana}{Una ventana tiene forma de rectángulo de 2 m de ancho por 1.5 m de alto, con un semicírculo encima (de diámetro 2 m). ¿Cuál es su área?}{Rectángulo: 3 m$^2$. Semicírculo de radio 1: $\tfrac{\pi}{2} \approx 1.57$ m$^2$. Total $\approx 4.57$ m$^2$.}
\ejx{La pizza}{Una pizza de 40 cm de diámetro se corta en 8 tajadas iguales. ¿Cuál es el área de una tajada?}{Área total $\pi \cdot 20^2 = 400\pi$. Una tajada: $50\pi \approx 157$ cm$^2$.}}
\masejemplosde{pu-ge-5}{
\ejx{El marco}{Un cuadro de $30 \times 20$ cm lleva un marco de 5 cm de ancho alrededor. ¿Cuál es el área del marco?}{Con marco: $40 \times 30 = 1\,200$. Sin marco: 600. El marco mide 600 cm$^2$.}
\ejx{Semicírculo en un rectángulo}{En un rectángulo de $10 \times 5$ se dibuja un semicírculo cuyo diámetro es el lado de 10. ¿Qué área queda fuera del semicírculo?}{$50 - \tfrac12 \pi \cdot 5^2 = 50 - 12.5\pi \approx 10.7$.}}
\masejemplosde{pu-ge-6}{
\ejx{La piscina}{Una piscina mide 8 m de largo, 4 m de ancho y 1.5 m de profundidad. ¿Cuántos litros caben?}{$8 \cdot 4 \cdot 1.5 = 48$ m$^3 = 48\,000$ litros.}
\ejx{Latas en una caja}{Una lata cilíndrica tiene 6 cm de diámetro y 10 cm de alto. ¿Cuántas caben, de pie, en una caja de $30 \times 24 \times 10$ cm?}{En la base caben $5 \times 4 = 20$ círculos de 6 cm de diámetro, y la altura alcanza para un solo piso: 20 latas.}}
% --- Libro 2 · Orden ---
\masejemplosde{pu-or-1}{
\ejx{El podio}{En una carrera de cinco, Sara llegó antes que Tito y después de Vera. Uma llegó de última. Willy llegó justo antes que Sara. ¿En qué puesto llegó Tito?}{Cadena $V < S < T$, con $W$ pegado antes de $S$ y $U$ en el 5: $V, W, S, T, U$. Tito llegó cuarto.}
\ejx{Estaturas}{Ana es más alta que Bea; Cata es más baja que Bea; Dora es más alta que Ana. ¿Quién es la segunda más alta?}{Orden de mayor a menor: Dora, Ana, Bea, Cata. La segunda es Ana.}}
\masejemplosde{pu-or-2}{
\ejx{Cuatro en la mesa}{Cuatro amigos en una mesa cuadrada, uno por lado: Inés frente a Joel; Karen no está a la derecha de Inés. ¿Quién está a la derecha de Inés?}{Los dos lados libres son para Karen y el cuarto, Luis. Como Karen no está a la derecha de Inés, es Luis.}
\ejx{Ocho puestos}{En una mesa redonda de 8 puestos, Mario está frente a Nora y Olga está dos puestos a la derecha de Mario. ¿Cuántos puestos separan a Olga de Nora por el lado más corto?}{Mario en el 1, Nora en el 5, Olga en el 3 (o en el 7). Entre el 3 y el 5 hay un puesto en medio: están a 2 puestos de distancia.}}
\masejemplosde{pu-or-3}{
\ejx{Tres oficios}{Paula, Quique y Rosa son chef, piloto y veterinaria. Paula no cocina. Quique le tiene miedo a volar. Rosa no trabaja con animales ni cocina. ¿Qué es cada uno?}{Rosa: piloto. Quique no es piloto, y Paula no es chef: Quique es chef y Paula veterinaria.}
\ejx{Las mascotas}{Tres niños tienen un gato, un perro y un pez. Sami no tiene gato. Tina tiene el pez. ¿Qué tiene Ugo?}{Tina, el pez. Sami no tiene gato: tiene el perro. Ugo, el gato.}}
\masejemplosde{pu-or-4}{
\ejx{Turnos}{Tres vigilantes (X, Y, Z) cubren lunes, martes y miércoles, uno por día. X no trabaja el lunes. Si Y trabaja el martes, entonces Z trabaja el lunes. ¿Cuántas asignaciones son posibles?}{El lunes es de Y o de Z. Si Y lunes: (Y, X, Z) y (Y, Z, X). Si Z lunes: (Z, X, Y) y (Z, Y, X). Las cuatro cumplen: 4.}
\ejx{Necesariamente}{Con las mismas reglas, ¿es necesariamente cierto que X no trabaja el lunes y que Z trabaja antes que X?}{Lo primero sí (es una regla). Lo segundo no: en (Y, X, Z), X trabaja antes que Z.}}
% --- Lógica ---
\masejemplosde{pu-lo-1}{
\ejx{Tres variables}{Si $p$ es V, $q$ es F y $r$ es F, evalúa $(p \wedge \neg q) \to (q \vee r)$.}{$p \wedge \neg q$: V. $q \vee r$: F. V $\to$ F es F.}
\ejx{Cuándo es verdadera}{¿Para qué valores de $p$ y $q$ es verdadera $\neg p \wedge (p \vee q)$?}{$\neg p$ exige $p$ F; entonces $p \vee q$ exige $q$ V. Solo con $p$ F y $q$ V.}}
\masejemplosde{pu-lo-2}{
\ejx{Negar un «o»}{Niega: «El sábado voy a cine o me quedo estudiando».}{«El sábado no voy a cine y no me quedo estudiando».}
\ejx{Negar con «algunos»}{Niega: «Algunos estudiantes no entregaron la tarea».}{«Todos los estudiantes entregaron la tarea».}}
\masejemplosde{pu-lo-3}{
\ejx{¿Qué se concluye?}{«Si un número termina en 0, es múltiplo de 5». El número $n$ es múltiplo de 5. ¿Termina en 0?}{No se puede asegurar (afirmar el consecuente): 15 es múltiplo de 5 y no termina en 0.}
\ejx{Cadena}{«Si apruebo el examen, entro a la universidad. Si entro a la universidad, me mudo a Medellín». No me mudé a Medellín. ¿Qué se concluye?}{Por tollens, no entré a la universidad; y otra vez por tollens, no aprobé el examen.}}
\masejemplosde{pu-lo-4}{
\ejx{Un silogismo válido}{«Todos los colibríes son aves. Ningún ave es mamífero». ¿Qué se concluye?}{El círculo de colibríes está dentro de aves, que no toca a mamíferos: ningún colibrí es mamífero.}
\ejx{Uno que no lo es}{«Todos los médicos estudiaron biología. Clara estudió biología». ¿Es médica Clara?}{No se sigue: Clara está en el círculo grande, pero puede estar fuera del pequeño.}}
\masejemplosde{pu-lo-5}{
\ejx{Dos hermanos}{Uno de dos hermanos siempre miente y el otro siempre dice la verdad. El primero dice: «Mi hermano es el mentiroso». ¿Se puede saber quién es quién?}{No: si el primero es veraz, la frase es cierta; si es mentiroso, la frase es falsa, y también cuadra. Las dos situaciones son posibles.}
\ejx{Tres cofres}{Hay tres cofres y solo uno tiene el tesoro. Rótulos: cofre 1, «El tesoro está aquí»; cofre 2, «El tesoro no está aquí»; cofre 3, «El tesoro no está en el cofre 1». Solo un rótulo dice la verdad. ¿Dónde está?}{Los rótulos 1 y 3 se contradicen: uno de ellos es el verdadero, así que el 2 miente: el tesoro está en el cofre 2. (Entonces el 1 miente y el 3 dice la verdad.)}}
% --- Conjuntos ---
\masejemplosde{pu-cj-1}{
\ejx{Operaciones encadenadas}{Con $U = \{1, \ldots, 9\}$, $A = \{1, 3, 5, 7, 9\}$ y $B = \{3, 6, 9\}$, halla $A - B$ y $(A \cap B)'$.}{$A - B = \{1, 5, 7\}$. $A \cap B = \{3, 9\}$, y su complemento es $\{1, 2, 4, 5, 6, 7, 8\}$.}
\ejx{Contar con la fórmula}{En un curso, 18 estudiantes tienen perro, 12 tienen gato y 5 tienen ambos. ¿Cuántos tienen al menos una de las dos mascotas?}{$18 + 12 - 5 = 25$.}}
\masejemplosde{pu-cj-2}{
\ejx{Tres idiomas}{De 60 personas: 30 hablan inglés, 25 francés y 20 alemán; 10 inglés y francés, 8 inglés y alemán, 6 francés y alemán; 4 los tres. ¿Cuántas no hablan ninguno?}{$30 + 25 + 20 - 10 - 8 - 6 + 4 = 55$ hablan alguno. Ninguno: 5.}
\ejx{Solo uno}{Con los mismos datos, ¿cuántas hablan solo inglés?}{$30 - (10 - 4) - (8 - 4) - 4 = 30 - 6 - 4 - 4 = 16$.}}
% --- Conteo ---
\masejemplosde{pu-co-1}{
\ejx{Claves con letras y números}{Una clave tiene 2 vocales seguidas de 2 dígitos. ¿Cuántas claves hay si nada se repite?}{$5 \cdot 4 \cdot 10 \cdot 9 = 1\,800$.}
\ejx{Rutas con escala}{De la ciudad A a la B hay 4 vuelos y de B a C hay 3 buses y 2 trenes. ¿De cuántas maneras se va de A a C pasando por B?}{$4 \cdot (3 + 2) = 20$.}}
\masejemplosde{pu-co-2}{
\ejx{Vocales juntas}{¿De cuántas maneras se ordenan las letras de MURCIELAGO si las cinco vocales deben quedar juntas?}{Bloque de vocales más 5 consonantes: $6!$ órdenes; dentro del bloque, $5!$. Total $720 \cdot 120 = 86\,400$.}
\ejx{Una foto}{Cuatro hombres y tres mujeres posan en fila, con las mujeres en los tres puestos del centro. ¿De cuántas maneras?}{$3!$ para las mujeres y $4!$ para los hombres: $6 \cdot 24 = 144$.}}
\masejemplosde{pu-co-3}{
\ejx{El examen}{En un examen hay que responder 8 de 10 preguntas. ¿De cuántas maneras se pueden elegir?}{$\binom{10}{8} = \binom{10}{2} = 45$.}
\ejx{Con una obligatoria}{Y si las tres primeras son obligatorias, ¿de cuántas maneras?}{Faltan 5 de las otras 7: $\binom75 = 21$.}}
\masejemplosde{pu-co-4}{
\ejx{Al menos una mujer}{De 5 hombres y 4 mujeres se elige un comité de 3. ¿Cuántos comités tienen al menos una mujer?}{Total $\binom93 = 84$; solo hombres $\binom53 = 10$. Con al menos una mujer: 74.}
\ejx{Cifras repetidas}{¿Cuántos números de tres cifras tienen al menos dos cifras iguales?}{Total 900; con las tres distintas $9 \cdot 9 \cdot 8 = 648$. Con alguna repetida: 252.}}
\masejemplosde{pu-co-5}{
\ejx{Pasando por la panadería}{En una cuadrícula, la casa está en $(0, 0)$, la panadería en $(2, 1)$ y el colegio en $(4, 3)$. ¿Cuántos caminos cortos pasan por la panadería?}{Hasta la panadería: $\binom31 = 3$. De ahí al colegio (2 a la derecha, 2 arriba): $\binom42 = 6$. Total 18.}
\ejx{Rectángulos}{¿Cuántos rectángulos (de cualquier tamaño) hay en una cuadrícula de $3 \times 3$?}{$\binom42 \cdot \binom42 = 36$.}}
% --- Probabilidad ---
\masejemplosde{pu-pb-1}{
\ejx{Una carta}{De una baraja de 52 cartas se saca una. ¿Probabilidad de que sea un as o una carta de corazones?}{4 ases $+$ 13 corazones $-$ 1 as de corazones $= 16$. $P = \tfrac{16}{52} = \tfrac{4}{13}$.}
\ejx{Dos dados distintos}{Se lanzan dos dados. ¿Probabilidad de que los dos números sean distintos?}{Iguales: 6 casos. Distintos: 30 de 36: $\tfrac56$.}}
\masejemplosde{pu-pb-2}{
\ejx{Tres semáforos}{La probabilidad de encontrar en verde cada uno de tres semáforos independientes es $\tfrac12$. ¿Probabilidad de encontrar al menos uno en rojo?}{Todos en verde: $\tfrac18$. Al menos uno en rojo: $\tfrac78$.}
\ejx{Dos calcetines}{En un cajón hay 4 medias negras y 2 blancas. Se sacan dos al azar. ¿Probabilidad de que formen un par del mismo color?}{Negras: $\tfrac46 \cdot \tfrac35 = \tfrac{12}{30}$. Blancas: $\tfrac26 \cdot \tfrac15 = \tfrac{2}{30}$. Total $\tfrac{14}{30} = \tfrac{7}{15}$.}}
% --- Secuencias ---
\masejemplosde{pu-se-1}{
\ejx{Multiplicar y restar}{¿Qué número sigue en $\;4,\ 7,\ 13,\ 25,\ 49,\ \ldots$?}{Las diferencias se duplican: 3, 6, 12, 24; sigue $+48$: 97. (También: doble menos 1.)}
\ejx{El puesto 30}{¿Cuál es el término 30 de $\;5,\ 8,\ 11,\ 14,\ \ldots$?}{$5 + 29 \cdot 3 = 92$.}}
\masejemplosde{pu-se-2}{
\ejx{Letras y números}{¿Qué sigue en $\;$A1, C4, E9, G16, …?}{Letras de dos en dos (A, C, E, G, I) y cuadrados (1, 4, 9, 16, 25): I25.}
\ejx{Hacia atrás}{¿Qué letra sigue en $\;$Y, V, S, P, …? (alfabeto de 26 letras)}{25, 22, 19, 16: baja de a 3. Sigue 13: M.}}
\masejemplosde{pu-se-3}{
\ejx{El collar}{Un collar repite cuentas en el orden rojo, rojo, azul, verde. ¿De qué color es la cuenta 150? ¿Cuántas rojas hay entre las primeras 150?}{$150 = 37 \cdot 4 + 2$: la segunda del ciclo, roja. Rojas: $37 \cdot 2 + 2 = 76$.}
\ejx{Medicamentos}{Un paciente toma una pastilla cada 8 horas y otra cada 12. Si toma las dos a las 6:00 a.\,m. del lunes, ¿cuándo vuelve a tomarlas juntas?}{MCM$(8, 12) = 24$ horas: el martes a las 6:00 a.\,m.}}
% --- Códigos ---
\masejemplosde{pu-cd-1}{
\ejx{Dos lugares}{Si PERRO se escribe RGTTQ, ¿cómo se escribe GATO?}{Cada letra avanza 2 lugares: G$\to$I, A$\to$C, T$\to$V, O$\to$Q. ICVQ.}
\ejx{Decodificar}{Con la misma clave, ¿qué palabra es ECUC?}{Se retroceden 2 lugares: C, A, S, A. CASA.}}
\masejemplosde{pu-cd-2}{
\ejx{Tres frases}{«zu ka» = «luz verde»; «ka mi» = «verde claro»; «ro zu» = «mucha luz». ¿Cómo se dice «mucho»?}{«ka» es común a las dos primeras, como «verde». Entonces zu = luz, y en la tercera, ro = mucha (mucho).}
\ejx{Parejas de números}{Si $(2, 3) \to 7$, $(4, 1) \to 17$ y $(5, 2) \to 27$, ¿a qué va $(3, 3)$?}{El resultado crece mucho con el primer número: se prueba su cuadrado. $2^2 + 3 = 7$, $4^2 + 1 = 17$, $5^2 + 2 = 27$: la regla es $a^2 + b$. Entonces $(3, 3) \to 12$.}}
% --- Espacial ---
\masejemplosde{pu-es-1}{
\ejx{El dron}{Un dron vuela 6 km al norte, 8 km al este y 6 km al sur. ¿A qué distancia en línea recta está de donde salió?}{Queda en $(8, 0)$: a 8 km al este.}
\ejx{En línea recta}{Una persona camina 9 cuadras al este y 12 al norte. ¿A cuántas cuadras está del inicio, en línea recta?}{Terna $9$-$12$-$15$: a 15 cuadras.}}
\masejemplosde{pu-es-2}{
\ejx{Caras pintadas}{Un cubo de $3 \times 3 \times 3$ formado por 27 cubitos se pinta por fuera. ¿Cuántos cubitos tienen exactamente una cara pintada?}{Los del centro de cada cara: 6.}
\ejx{La torre}{Se apilan cubos formando una escalera de 4 escalones (4, 3, 2 y 1 cubos de alto, uno de fondo). ¿Cuántos cubos tiene?}{$4 + 3 + 2 + 1 = 10$.}}
\masejemplosde{pu-es-3}{
\ejx{El dado}{En un dado normal se ven el 1, el 2 y el 3. ¿Cuánto suman las caras que no se ven?}{$21 - 6 = 15$ (son el 6, el 5 y el 4).}
\ejx{La fila de cuatro}{En un desarrollo, cuatro caras están en fila (P, Q, R, S), con T encima de Q y U debajo de Q. ¿Cuál es opuesta a P?}{En la fila, la que está a una de distancia: R. (Q es opuesta a S, y T a U.)}}
\masejemplosde{pu-es-4}{
\ejx{La hora del espejo}{Un reloj de agujas en el espejo parece marcar las 5:40. ¿Qué hora es?}{$12{:}00 - 5{:}40 = 6{:}20$.}
\ejx{Doblar y cortar}{Una hoja se dobla por la mitad dos veces y se corta una esquina que no toca ningún doblez. ¿Cuántos cortes aparecen al abrirla?}{Cada doblez duplica: $2^2 = 4$ cortes, uno en cada esquina de la hoja.}}
% --- Gráficas ---
\masejemplosde{pu-gr-1}{
\ejx{La torta del presupuesto}{En una torta de gastos, arriendo ocupa $126^\circ$. ¿Qué porcentaje es?}{$\tfrac{126}{360} = 0.35$: el $35\,\%$.}
\ejx{Variación porcentual}{Las visitas a un museo pasaron de 1\,250 a 1\,000. ¿En qué porcentaje bajaron?}{$\tfrac{250}{1\,250} = 20\,\%$.}}
\masejemplosde{pu-gr-2}{
\ejx{El dato que falta}{Los datos 4, 7, 9, $x$ y 10 tienen media 8. ¿Cuánto vale $x$ y cuál es la mediana?}{Suma 40: $x = 10$. Ordenados: 4, 7, 9, 10, 10. Mediana 9.}
\ejx{Dos modas}{Halla la moda de 3, 5, 5, 6, 8, 8, 9.}{El 5 y el 8 aparecen dos veces: hay dos modas (es bimodal).}}
\masejemplosde{pu-gr-3}{
\ejx{Lo que sí se puede decir}{Una tabla muestra que el colegio A tuvo 120 graduados y el B 80. ¿Se puede concluir que el A es mejor colegio?}{No: solo se puede decir que el A graduó a 40 estudiantes más. «Mejor» es un juicio, y además no se sabe cuántos estudiantes tiene cada uno.}
\ejx{Porcentajes de totales distintos}{En el colegio A aprobó el $80\,\%$ de 50 estudiantes y en el B el $60\,\%$ de 100. ¿Dónde aprobaron más estudiantes?}{A: 40. B: 60. Aprobaron más en el B, aunque su porcentaje sea menor.}}
